\documentclass[12pt,a4paper]{report}

% ============================================================
% PACKAGES
% ============================================================

\usepackage[utf8]{inputenc}
\usepackage[T1]{fontenc}
\usepackage[french]{babel}

\usepackage{graphicx}
\usepackage{float}
\usepackage{geometry}
\usepackage{setspace}
\usepackage{amsmath}
\usepackage{amssymb}
\usepackage{booktabs}
\usepackage{array}
\usepackage{tabularx}
\usepackage{longtable}
\usepackage{enumitem}
\usepackage{xcolor}
\usepackage{listings}
\usepackage{url}
\usepackage{hyperref}

% ============================================================
% MISE EN PAGE
% ============================================================

\geometry{
    left=2.5cm,
    right=2.5cm,
    top=2.5cm,
    bottom=2.5cm
}

\onehalfspacing

\setlength{\parindent}{0.7cm}
\setlength{\parskip}{0.2cm}

\renewcommand{\arraystretch}{1.4}

% ============================================================
% HYPERREF
% ============================================================

\hypersetup{
    colorlinks=true,
    linkcolor=black,
    urlcolor=blue,
    citecolor=black
}

% ============================================================
% CODE PYTHON
% ============================================================

\lstset{
    language=Python,
    basicstyle=\ttfamily\small,
    breaklines=true,
    frame=single,
    numbers=left,
    numberstyle=\tiny,
    showstringspaces=false,
    columns=fullflexible,
    keepspaces=true
}

% ============================================================
% DOCUMENT
% ============================================================

\begin{document}

% ============================================================
% PAGE DE GARDE
% ============================================================

\begin{titlepage}
\centering

\vspace*{0.5cm}

{\Large \textbf{Université de Kinshasa}}\\[0.2cm]
{\large Faculté des Sciences et Technologies}\\[0.2cm]
{\large Mention : Mathématiques, Statistique et Informatique}\\[0.3cm]

{\large Cours : Science des données}\\[0.2cm]
{\large Master 1, Semestre 2, 2025--2026}\\[0.5cm]

% --- Logo ---
\includegraphics[width=3.5cm]{unikin.PNG}\\[0.5cm]

% --- Titre ---
{\Large \textbf{Travail Pratique de Science des données}}\\[0.5cm]

\rule{\linewidth}{0.5mm}\\[0.3cm]

{\large \bfseries
IMPLÉMENTATION DE L'ALGORITHME DES NUÉES DYNAMIQUES\\
(From Scratch)
}\\[0.5cm]

\rule{\linewidth}{0.5mm}\\[0.8cm]

% --- Auteur ---
\raggedright
\textbf{Présenté par :}\\[0.2cm]

ZOLA NDIENE Priscilia

\vfill

% --- Enseignant ---
\centering

{\large \textbf{Professeur : KAFUNDA}}\\[0.2cm]
Doct. KAMINGU Gradi\\[0.5cm]

{\large \textbf{Année académique 2025--2026}}

\end{titlepage}


% ============================================================
% TABLE DES MATIÈRES
% ============================================================

\tableofcontents

\newpage

% ============================================================
% RÉSUMÉ
% ============================================================

\chapter*{Résumé}
\addcontentsline{toc}{chapter}{Résumé}

Ce travail porte sur la conception et l'implémentation d'une
application des Nuées dynamiques permettant d'utiliser plusieurs
types de représentations des classes.

Contrairement à une implémentation limitée au modèle K-means,
l'application proposée considère les Nuées dynamiques comme un cadre
général dans lequel une classe peut être représentée par un point,
un ensemble de points, des axes ou composantes factorielles, une
distribution de probabilités ou encore une structure représentative.

L'utilisateur peut choisir le nombre de classes ainsi que le type
de représentation avant l'exécution de l'algorithme. Les paramètres
spécifiques sont ensuite adaptés à la représentation choisie.

L'algorithme général repose sur quatre éléments fondamentaux :
la représentation de la classe, le critère d'affectation, la mise à
jour de la représentation et le critère de convergence.

L'implémentation est réalisée en Python à partir de zéro, sans utiliser
une fonction de clustering préexistante telle que \texttt{KMeans()} de
Scikit-learn.

\textbf{Mots-clés :} Nuées dynamiques, Diday, clustering,
K-means, représentation, prototypes, axes, distribution,
Python, Machine Learning.

% ============================================================
% INTRODUCTION
% ============================================================

\chapter*{Introduction générale}
\addcontentsline{toc}{chapter}{Introduction générale}

Le clustering constitue une importante famille de méthodes
d'apprentissage non supervisé. Il permet de regrouper des individus
présentant des caractéristiques similaires sans disposer nécessairement
de classes connues à l'avance.

Parmi les méthodes de classification automatique figurent les
Nuées dynamiques. Cette approche ne se limite pas à un seul type de
représentation d'une classe. Elle permet de définir une représentation
adaptée à la nature des données et au problème étudié.

Une représentation particulièrement connue est le point représentatif,
généralement le centroïde. Dans ce cas particulier, les Nuées dynamiques
conduisent au principe du K-means.

Cependant, limiter les Nuées dynamiques aux seuls centroïdes ne permet
pas de mettre en évidence leur caractère général. Une classe peut
également être représentée par plusieurs points, des axes ou composantes
factorielles, une distribution de probabilités ou une structure
représentative.

L'objectif de ce travail est donc de concevoir une application permettant
à l'utilisateur de choisir le type de représentation avant d'exécuter
l'algorithme.

% ============================================================
% PROBLÉMATIQUE
% ============================================================

\section*{Problématique}
\addcontentsline{toc}{section}{Problématique}

Comment concevoir et implémenter une application des Nuées dynamiques
capable de gérer plusieurs types de représentations des classes, tout
en permettant d'adapter le critère d'affectation, la mise à jour et
le critère de convergence à la représentation choisie ?

% ============================================================
% OBJECTIFS
% ============================================================

\section*{Objectif général}
\addcontentsline{toc}{section}{Objectif général}

L'objectif général est de concevoir et d'implémenter une application
générale des Nuées dynamiques permettant de choisir différentes
représentations des classes.

\section*{Objectifs spécifiques}
\addcontentsline{toc}{section}{Objectifs spécifiques}

Les objectifs spécifiques sont :

\begin{itemize}
    \item comprendre le principe général des Nuées dynamiques ;
    \item distinguer les Nuées dynamiques du cas particulier des
    K-means ;
    \item permettre le choix du nombre de classes ;
    \item permettre le choix du type de représentation ;
    \item adapter le critère d'affectation à la représentation choisie ;
    \item mettre à jour les représentations des classes ;
    \item définir un critère de convergence ;
    \item visualiser les résultats obtenus ;
    \item implémenter les mécanismes principaux sans utiliser
    une fonction K-means existante.
\end{itemize}

% ============================================================
% CHAPITRE 1
% ============================================================

\chapter{Fondements théoriques}

\section{Classification non supervisée}

La classification non supervisée consiste à rechercher automatiquement
des groupes dans un ensemble de données sans disposer de classes
préalablement étiquetées.

Le clustering constitue l'une des principales approches de classification
non supervisée.

\section{Principe des Nuées dynamiques}

Les Nuées dynamiques constituent un cadre général de classification
dans lequel chaque classe est décrite par une représentation appelée
nuée.

Le fonctionnement général peut être résumé par les étapes suivantes :

\begin{enumerate}
    \item choisir le nombre de classes ;
    \item choisir une représentation des classes ;
    \item initialiser les représentations ;
    \item affecter les individus aux classes ;
    \item mettre à jour les représentations ;
    \item calculer un critère d'optimisation ;
    \item vérifier la convergence ;
    \item répéter les étapes jusqu'à convergence.
\end{enumerate}

\section{Représentation d'une classe}

La représentation d'une classe dépend de la nature des données et du
problème étudié.

Dans cette application, cinq représentations sont considérées.

\begin{table}[H]
\centering
\caption{Types de représentations considérés}
\label{tab:representations}
\begin{tabular}{|c|l|p{8cm}|}
\hline
\textbf{N°} & \textbf{Représentation} & \textbf{Principe} \\
\hline
1 & Point représentatif &
Une classe est représentée par un seul point, généralement le centroïde. \\
\hline
2 & Ensemble de points représentatifs &
Une classe est représentée par plusieurs points prototypes. \\
\hline
3 & Axes / composantes factorielles &
Une classe est représentée par un centre et une structure multidimensionnelle fondée sur des axes principaux. \\
\hline
4 & Distribution de probabilités &
Une classe est décrite par une distribution caractérisée ici par une moyenne et une covariance. \\
\hline
5 & Structure représentative &
Une classe est décrite par une structure adaptée aux données ; l'implémentation proposée utilise une structure d'intervalle multidimensionnel. \\
\hline
\end{tabular}
\end{table}

\section{Point représentatif et cas particulier des K-means}

Lorsque chaque classe est représentée par un seul point, le point
représentatif peut être choisi comme le centroïde des observations
de la classe.

La distance euclidienne peut alors être utilisée pour affecter chaque
observation au centroïde le plus proche.

Cette configuration correspond au cas classique des K-means.

Ainsi :

\[
\boxed{
\text{K-means}
\subset
\text{Nuées dynamiques}
}
\]

dans le cadre de cette modélisation.

\section{Ensemble de points représentatifs}

Une classe peut également être représentée par plusieurs points
représentatifs.

Si une classe contient plusieurs sous-structures, un seul centroïde
peut ne pas suffire à décrire correctement sa forme.

L'utilisation de plusieurs prototypes permet alors de représenter
plusieurs zones de la classe.

Dans l'application, la distance d'un individu à une classe est définie
comme la distance au représentant le plus proche.

\section{Axes ou composantes factorielles}

Une classe peut être représentée par une structure multidimensionnelle
composée d'un centre et d'un ou plusieurs axes.

Dans l'implémentation proposée, les axes sont déterminés à partir de
la matrice de covariance des observations de la classe.

La distance utilisée correspond à la distance entre l'individu et
l'espace engendré par les axes de représentation.

\section{Distribution de probabilités}

Une classe peut être représentée par une distribution statistique.

Dans l'application, une représentation gaussienne simple est utilisée.
Elle est définie par :

\[
R_j = (\mu_j,\Sigma_j)
\]

où $\mu_j$ représente la moyenne de la classe et $\Sigma_j$ sa matrice
de covariance.

La distance de Mahalanobis est utilisée pour comparer une observation
à la distribution.

\[
d_M(x,\mu_j)
=
\sqrt{
(x-\mu_j)^T
\Sigma_j^{-1}
(x-\mu_j)
}
\]

\section{Structure représentative}

La représentation structurelle est prévue comme une représentation
plus générale et extensible.

Dans cette première implémentation, une classe est représentée par
un intervalle multidimensionnel défini par un minimum et un maximum
pour chaque variable.

Cette représentation permet de décrire une région occupée par les
observations de la classe.

% ============================================================
% CHAPITRE 2
% ============================================================

\chapter{Conception et implémentation}

\section{Architecture générale}

Le fonctionnement de l'application est organisé comme suit :

\begin{verbatim}
Données
   |
   v
Choix du nombre de classes K
   |
   v
Choix de la représentation
   |
   v
Initialisation des représentations
   |
   v
Affectation des individus
   |
   v
Mise à jour des représentations
   |
   v
Calcul du critère
   |
   v
Test de convergence
   |
   +------ Non ------> Nouvelle itération
   |
   Oui
   |
   v
Résultats
\end{verbatim}

\section{Choix de la représentation}

Au démarrage, l'utilisateur peut sélectionner l'une des cinq
représentations.

\begin{table}[H]
\centering
\caption{Menu de choix de l'application}
\label{tab:menu}
\begin{tabular}{|c|l|l|}
\hline
\textbf{Choix} & \textbf{Représentation} & \textbf{Paramètre spécifique} \\
\hline
1 & Point représentatif & Type de point : centroïde \\
\hline
2 & Ensemble de points & Nombre de représentants par classe \\
\hline
3 & Axes / composantes & Nombre d'axes \\
\hline
4 & Distribution de probabilités & Moyenne + covariance \\
\hline
5 & Structure représentative & Structure min/max \\
\hline
\end{tabular}
\end{table}

\section{Paramètres généraux}

Certains paramètres sont communs à toutes les représentations.

\begin{table}[H]
\centering
\caption{Paramètres généraux de l'application}
\label{tab:parametres}
\begin{tabular}{|l|c|p{7cm}|}
\hline
\textbf{Paramètre} & \textbf{Valeur} & \textbf{Description} \\
\hline
Nombre de classes $K$ & 3 & Nombre de classes recherchées dans l'expérimentation. \\
\hline
Nombre maximal d'itérations & 100 & Nombre maximal d'itérations autorisées. \\
\hline
Tolérance & $10^{-4}$ & Seuil utilisé pour la convergence. \\
\hline
Graine aléatoire & 42 & Permet de reproduire l'initialisation. \\
\hline
Nombre d'observations & 90 & Taille du jeu de données artificiel. \\
\hline
Dimension & 2 & Nombre de variables utilisées. \\
\hline
\end{tabular}
\end{table}

\section{Les quatre éléments fondamentaux}

L'algorithme est organisé autour de quatre éléments.

\begin{table}[H]
\centering
\caption{Éléments fondamentaux des Nuées dynamiques}
\label{tab:elements}
\begin{tabular}{|c|l|p{8cm}|}
\hline
\textbf{N°} & \textbf{Élément} & \textbf{Rôle} \\
\hline
1 & Représentation &
Décrit la manière dont chaque classe est représentée. \\
\hline
2 & Critère d'affectation &
Détermine la classe à laquelle un individu doit être affecté. \\
\hline
3 & Mise à jour &
Construit une nouvelle représentation à partir des individus affectés. \\
\hline
4 & Convergence &
Détermine si les représentations sont suffisamment stables pour arrêter l'algorithme. \\
\hline
\end{tabular}
\end{table}

\section{Correspondance entre représentations et critères}

\begin{table}[H]
\centering
\caption{Critères utilisés selon la représentation}
\label{tab:criteres}
\begin{tabular}{|l|l|l|}
\hline
\textbf{Représentation} & \textbf{Critère} & \textbf{Mise à jour} \\
\hline
Point & Distance euclidienne & Moyenne des observations \\
\hline
Ensemble de points & Minimum des distances & Moyenne des sous-groupes \\
\hline
Axes & Distance à l'espace représenté & Centre + axes principaux \\
\hline
Distribution & Distance de Mahalanobis & Moyenne + covariance \\
\hline
Structure & Distance à la structure & Minimum + maximum \\
\hline
\end{tabular}
\end{table}

\section{Jeu de données}

L'expérimentation utilise un jeu de données artificiel constitué de
90 observations en deux dimensions.

\begin{table}[H]
\centering
\caption{Caractéristiques des données}
\label{tab:donnees}
\begin{tabular}{|l|c|p{7cm}|}
\hline
\textbf{Caractéristique} & \textbf{Valeur} & \textbf{Description} \\
\hline
Nombre d'observations & 90 & Points générés artificiellement. \\
\hline
Nombre de variables & 2 & Coordonnées utilisées pour la représentation. \\
\hline
Nombre de groupes recherchés & 3 & Valeur utilisée pour l'expérimentation. \\
\hline
Observations par groupe initial & 30 & Répartition utilisée lors de la génération. \\
\hline
\end{tabular}
\end{table}

\section{Organisation du programme}

\begin{table}[H]
\centering
\caption{Principales fonctions du programme}
\label{tab:fonctions}
\begin{tabular}{|l|p{8.5cm}|}
\hline
\textbf{Fonction} & \textbf{Rôle} \\
\hline
\texttt{generer\_donnees()} &
Génère le jeu de données artificiel. \\
\hline
\texttt{initialiser()} &
Initialise les représentations selon le choix de l'utilisateur. \\
\hline
\texttt{affecter()} &
Affecte les observations aux classes selon le critère choisi. \\
\hline
\texttt{mettre\_a\_jour()} &
Met à jour les représentations des classes. \\
\hline
\texttt{calculer\_critere()} &
Calcule le critère d'optimisation. \\
\hline
\texttt{distance\_representations()} &
Mesure le déplacement des représentations entre deux itérations. \\
\hline
\texttt{nuees\_dynamiques()} &
Coordonne l'ensemble des étapes de l'algorithme. \\
\hline
\texttt{afficher\_resultats()} &
Produit les visualisations finales. \\
\hline
\end{tabular}
\end{table}

% ============================================================
% CHAPITRE 3
% ============================================================

\chapter{Expérimentation et résultats}

\section{Déroulement de l'expérimentation}

L'utilisateur commence par sélectionner le type de représentation.
Il indique ensuite le nombre de classes $K$, puis les paramètres
spécifiques à la représentation sélectionnée.

L'algorithme initialise les représentations et commence les itérations.

À chaque itération :

\begin{enumerate}
    \item les observations sont affectées aux classes ;
    \item les représentations sont mises à jour ;
    \item le critère d'optimisation est calculé ;
    \item le déplacement des représentations est calculé ;
    \item la condition de convergence est vérifiée.
\end{enumerate}

\section{Résultats graphiques}

L'application génère une représentation graphique des classes.

\begin{figure}[H]
    \centering
    \includegraphics[width=0.85\textwidth]{Cléo.PNG}
    \caption{Visualisation des classes obtenues}
    \label{fig:clusters}
\end{figure}

L'application produit également une courbe montrant l'évolution
du critère d'optimisation.

\begin{figure}[H]
    \centering
    \includegraphics[width=0.85\textwidth]{cléa.PNG}
    \caption{Évolution du critère d'optimisation}
    \label{fig:critere}
\end{figure}

\section{Analyse des résultats}

Les résultats permettent d'observer l'influence du choix de la
représentation sur le processus de classification.

Dans le cas du point représentatif, chaque classe est résumée par
un seul centroïde. Cette configuration correspond au cas particulier
du K-means.

Dans le cas de l'ensemble de points représentatifs, plusieurs
prototypes sont utilisés pour décrire une même classe. Cette
représentation permet de tenir compte de plusieurs zones au sein
d'une classe.

La représentation par axes permet de tenir compte d'une organisation
multidimensionnelle des observations.

La représentation probabiliste prend en compte la moyenne et la
dispersion des observations à travers la matrice de covariance.

Enfin, la structure représentative fournit une représentation
extensible pouvant être adaptée à d'autres types de données ou
d'autres critères.

\section{Comparaison conceptuelle}

\begin{table}[H]
\centering
\caption{Comparaison des représentations}
\label{tab:comparaison}
\begin{tabular}{|l|p{3.5cm}|p{4cm}|p{3cm}|}
\hline
\textbf{Représentation} &
\textbf{Éléments représentatifs} &
\textbf{Information décrite} &
\textbf{Exemple de critère} \\
\hline
Point &
Un point &
Position centrale &
Euclidien \\
\hline
Ensemble de points &
Plusieurs points &
Plusieurs zones de la classe &
Minimum des distances \\
\hline
Axes &
Centre + axes &
Organisation multidimensionnelle &
Distance à l'espace \\
\hline
Distribution &
Moyenne + covariance &
Position et dispersion &
Mahalanobis \\
\hline
Structure &
Bornes min/max &
Région occupée &
Distance à la structure \\
\hline
\end{tabular}
\end{table}

\section{Avantages}

L'application présente plusieurs avantages :

\begin{itemize}
    \item elle distingue clairement les Nuées dynamiques du seul
    K-means ;
    \item elle permet de sélectionner une représentation ;
    \item les critères sont adaptés aux représentations ;
    \item l'architecture peut être étendue à d'autres représentations ;
    \item l'algorithme principal est implémenté manuellement ;
    \item les résultats peuvent être visualisés.
\end{itemize}

\section{Limites}

Cette première version présente certaines limites :

\begin{itemize}
    \item les données utilisées pour l'expérimentation sont artificielles ;
    \item les représentations proposées ne couvrent pas toutes les
    possibilités théoriques des Nuées dynamiques ;
    \item la représentation structurelle utilisée constitue une
    première structure générique ;
    \item certaines représentations nécessitent des critères plus
    spécialisés selon la nature des données ;
    \item l'application est principalement destinée à l'étude et à
    la compréhension de l'algorithme.
\end{itemize}

% ============================================================
% CONCLUSION
% ============================================================

\chapter*{Conclusion générale}
\addcontentsline{toc}{chapter}{Conclusion générale}

Ce travail a permis de concevoir une application des Nuées dynamiques
en considérant la représentation d'une classe comme un élément central
de l'algorithme.

Contrairement à une implémentation limitée au K-means, l'application
permet de choisir entre plusieurs représentations : point représentatif,
ensemble de points représentatifs, axes ou composantes factorielles,
distribution de probabilités et structure représentative.

Cette conception permet de mettre en évidence le fait que les K-means
constituent un cas particulier correspondant à une représentation de
chaque classe par un seul point, généralement le centroïde.

L'algorithme implémenté repose sur quatre éléments essentiels :
la représentation, le critère d'affectation, la mise à jour de la
représentation et le critère de convergence.

L'architecture développée reste également extensible, puisqu'une
nouvelle représentation peut être ajoutée en définissant sa méthode
d'initialisation, son critère d'affectation, sa méthode de mise à jour
et son critère de convergence.

% ============================================================
% RÉFÉRENCES
% ============================================================

\chapter*{Références}
\addcontentsline{toc}{chapter}{Références}

\begin{enumerate}

    \item Diday, E. et Simon, J. C.
    Travaux sur les méthodes des Nuées dynamiques et la classification
    automatique.

    \item Bishop, C. M. (2006).
    \textit{Pattern Recognition and Machine Learning}.
    Springer.

    \item Han, J., Kamber, M. et Pei, J. (2012).
    \textit{Data Mining: Concepts and Techniques}.
    Morgan Kaufmann.

    \item Documentation Python.

    \item Documentation NumPy.

    \item Documentation Matplotlib.

\end{enumerate}

% ============================================================
% ANNEXE
% ============================================================

\appendix

\chapter{Extrait du code Python}

Le programme complet est disponible dans le dépôt GitHub associé au
travail. L'extrait suivant présente la fonction principale qui orchestre
les quatre éléments fondamentaux des Nuées dynamiques.

\begin{lstlisting}
def nuees_dynamiques(
    X,
    k,
    type_representation,
    parametre=None,
    max_iter=100,
    tolerance=1e-4,
    seed=42
):

    representations = initialiser(
        X,
        k,
        type_representation,
        parametre,
        seed
    )

    historique_critere = []

    for iteration in range(
        1,
        max_iter + 1
    ):

        etiquettes = affecter(
            X,
            representations,
            type_representation
        )

        nouvelles_representations = mettre_a_jour(
            X,
            etiquettes,
            representations,
            type_representation,
            parametre
        )

        critere = calculer_critere(
            X,
            etiquettes,
            nouvelles_representations,
            type_representation
        )

        historique_critere.append(
            critere
        )

        deplacement = distance_representations(
            representations,
            nouvelles_representations,
            type_representation
        )

        representations = (
            nouvelles_representations
        )

        if deplacement < tolerance:
            break

    return (
        etiquettes,
        representations,
        historique_critere,
        iteration
    )
\end{lstlisting}

\end{document}